\documentclass[12pt,a4j]{jarticle}
\usepackage{graphicx}
\begin{document}
\title{基礎プログラミングおよび演習 レポート ＃06}
\author{学籍番号, 氏名 (ペア: 氏名・学籍番号または「個人作業」)}
\date{提出日付}
\maketitle

\section{構想・計画・設計}

(どのような構想で絵を生成したか、具体的にどのように計画し、プログラムはどう設計したか)

\section{プログラムコード}
\begin{verbatim}
(ここにプログラムのソースコードを入れる)
\end{verbatim}
(必ず動作するものを提出してください。また絵を生成するために呼び出すRuby命令を最後の行に追加する。)
(「ruby ファイル名」で実行できるようにするため。)

\section{プログラムの説明}

(プログラムのどの部分が何をしているかを説明する)

\section{生成された絵}

(どんな絵という説明を書くこと)
(絵の画像ファイル形式はPostScriptであること。プログラムと絵が一致してい
ること。ファイル名は適宜変更してよい。)
\begin{center}
\includegraphics[width=12cm]{mypicture.ps}
\end{center}


\section{考察}

(考察は必須かつ重要。課題をやって分かったことや感想など。)

\section{アンケート}

\subsection{Q1：画像が自由に生成できるようになりましたか。}

(ここにQ1の回答を記入)

\subsection{Q2：画像をうまく生成する「コツ」は何だと思いましたか。}

(ここにQ2の回答を記入)

\subsection{Q3：リフレクション(今回の課題で分かったこと)・感想・要望をどうぞ。}

(ここにQ3の回答を記入)

\end{document}